\section{The finiteness program}

\begin{convention}[Arithmetic setting]
\label{conv:milne-finiteness-arithmetic}

\dcref{MI:ch4:1}
In this chapter $K$ is a number field, $S$ is a finite set of its places, and
dimensions and genera are fixed when a finiteness assertion says so.  Good or
semistable reduction is always understood at the places outside $S$.
\end{convention}

\begin{theorem}[Finiteness I]
\label{thm:milne-finiteness-I}

\dcref{MI:ch4:1.1}
For an abelian variety $A$ over a number field $K$, only finitely many
$K$-isomorphism classes of abelian varieties are isogenous to $A$.
\end{theorem}

\begin{proof}[Source proof]

When $A$ has semistable reduction everywhere, the Faltings-height theorem in
Section~IV.8 and modular-height finiteness give only finitely many isomorphism
classes in its isogeny class.  For general $A$, pass to a finite extension over
which it is semistable and apply the finite descent step for the forms over the
original number field.

\end{proof}


\begin{theorem}[Semisimplicity]
\label{thm:milne-finiteness-semisimplicity}

\dcref{MI:ch4:1.2}
Assume Finiteness I.  For every prime $\ell$, the representation of
$\Gal(K^{\mathrm{sep}}/K)$ on $V_\ell A$ is semisimple.
\end{theorem}

\begin{proof}[Source proof]

Under Finiteness~I, the Tate-module argument of Section~IV.2 realizes every
stable subspace by an endomorphism and obtains a stable complementary summand.

\end{proof}


\begin{theorem}[Tate homomorphism conjecture]
\label{thm:milne-finiteness-tate-conjecture}

\dcref{MI:ch4:1.3}
Assume Finiteness I for $A$, $B$, and $A\times B$.  For abelian varieties
$A,B/K$, the canonical map
$\Hom_K(A,B)\otimes\Z_\ell\to\Hom_{\Gal(K^{\mathrm{sep}}/K)}(T_\ell A,T_\ell B)$
is an isomorphism.
\end{theorem}

\begin{proof}[Source proof]

The proof is the rational Tate homomorphism theorem of Section~IV.2 together
with the integral finite-level torsion sequences.

\end{proof}


\begin{theorem}[Finiteness II]
\label{thm:milne-finiteness-II}

\dcref{MI:ch4:1.4}
For fixed $K$, dimension $g$, and finite set $S$ of finite places, only
finitely many $K$-isomorphism classes of $g$-dimensional abelian varieties have
good reduction outside $S$.
\end{theorem}

\begin{proof}[Source proof]

Section~IV.3 detects finitely many isogeny classes from Frobenius polynomials at
a finite set of places; Finiteness~I then makes each class finite.

\end{proof}


\begin{theorem}[Shafarevich finiteness]
\label{thm:milne-finiteness-shafarevich}

\dcref{MI:ch4:1.5}
For fixed $K$, genus $g$, and finite set $S$, only finitely many smooth complete
curves of genus $g$ have good reduction outside $S$.
\end{theorem}

\begin{proof}[Source proof]

Apply Finiteness~II to the principally polarized Jacobians and use the rational
Torelli theorem of Section~IV.4.

\end{proof}


\begin{theorem}[Faltings-height finiteness]
\label{thm:milne-finiteness-height}

\dcref{MI:ch4:1.6}
If $A/K$ has semistable reduction everywhere, the set of Faltings heights of
abelian varieties over $K$ isogenous to $A$ is finite.
\end{theorem}

\begin{proof}[Source proof]

This is the height theorem completed in Section~IV.8: isogenous varieties have
bounded stable Faltings height, and modular-height finiteness turns bounded
height into finitely many isomorphism classes.

\end{proof}


\section{Tate modules and semisimplicity}

\begin{lemma}[Finite torsion kernels of prime-to-characteristic isogenies]
\label{lem:milne-finite-level-tate}

\dcref{MI:ch4:2.1}
\uses{def:milne-tate-module,def:milne-isogeny}
Let $u:A\to B$ be an isogeny over a field $K$, with degree prime to
$\operatorname{char}(K)$.  The geometric kernel of $u$ is finite and Galois
stable, and every such finite stable subgroup occurs as the kernel of an
isogeny whose quotient is defined over $K$.
\end{lemma}

\begin{proof}

Over a separable closure this is the finite-kernel quotient construction for
abelian varieties.  If $N\subset A(K^{\mathrm{sep}})$ is stable under Galois,
the quotient $A/N$ and the quotient map descend to $K$.

\end{proof}


\begin{lemma}[Stable subspaces from endomorphisms]
\label{lem:milne-stable-tate-subspace}

\dcref{MI:ch4:2.4}
\uses{lem:milne-finite-level-tate,thm:milne-hom-tate-injective}
Assume that $A$ satisfies (*): up to isomorphism, only finitely many abelian
varieties are isogenous to $A$ by an isogeny of $\ell$-power degree.  For every
Galois-stable subspace $W\subset V_\ell A$ there is
$u\in\End(A)\otimes\Q_\ell$ with $uV_\ell A=W$.
\end{lemma}

\begin{proof}

Write $T=T_\ell A$ and $V=V_\ell A$.  Set
$X_n=(T\cap W)+\ell^nT$.  Each $X_n$ is Galois-stable and of finite index in
$T$, so Lemma~\ref{lem:milne-stable-tate-lattices} gives an isogeny
$f_n:B_n\to A$ with $f_n(T_\ell B_n)=X_n$.  By (*), infinitely many of the
$B_n$ are isomorphic; choose an infinite set $I$ and isomorphisms
$v_i:B_{i_0}\to B_i$ for a fixed $i_0\in I$.  Then
$u_i=f_i v_i f_{i_0}^{-1}$ lies in $\End(A)\otimes\Q_\ell$ and satisfies
$u_i(X_{i_0})=X_i$.

For $i\ge i_0$, $X_i\subset X_{i_0}$, so the $u_i$ lie in the compact set
$\Hom_{\mathbb Z_\ell}(X_{i_0},T)$.  After passing to a subsequence, they
converge to $u$.  The subspace $\End(A)\otimes\Q_\ell$ is closed in
$\End_{Q_\ell}(V)$, hence $u$ still belongs to it.  For $x\in X_{i_0}$,
$u(x)\in\bigcap_{i\in I}X_i=T\cap W$.  Conversely, if
$y\in\bigcap_{i\in I}X_i$, choose $x_i\in X_{i_0}$ with $u_i(x_i)=y$.
Compactness gives a convergent subsequence $x_i\to x\in X_{i_0}$, and then
$u(x)=y$.  Thus $u(X_{i_0})=T\cap W$, and after tensoring with $\Q_\ell$,
$uV=W$.

\end{proof}


\begin{corollary}[Tate's homomorphism theorem]
\label{cor:milne-tate-hom-finiteness}

\dcref{MI:ch4:2.6}
\uses{lem:milne-stable-tate-subspace,lem:milne-finite-level-tate}
Assuming Faltings finiteness for $A$, $B$, and $A\times B$,
\[
 \Hom(A,B)\otimes\Q_\ell
 \simeq \Hom_{\Gal(K^{\mathrm{sep}}/K)}(V_\ell A,V_\ell B).
\]
\end{corollary}

\begin{proof}[Source proof]

Apply the stable-subspace argument to the relevant products.  Equivalently,
compare the two inclusions in the matrix decomposition
\[
\End(A\times B)\otimes\Q_\ell
 =\End(A)\otimes\Q_\ell\oplus\Hom(B,A)\otimes\Q_\ell
 \oplus\Hom(A,B)\otimes\Q_\ell\oplus\End(B)\otimes\Q_\ell
\]
and its Galois-equivariant Tate-module analogue.  The endomorphism equality
for the product makes the middle summands equal, proving the asserted rational
isomorphism; the finite-level torsion sequences give the corresponding
integral injectivity.

\end{proof}


\begin{theorem}[Semisimplicity of Tate representations]
\label{thm:milne-tate-semisimplicity}

\dcref{MI:ch4:2.5}
\uses{cor:milne-tate-hom-finiteness}
Under the same finiteness hypotheses, $V_\ell A$ is a semisimple representation
of the absolute Galois group, and the canonical map
\[
\End(A)\otimes\Q_\ell\longrightarrow
\End_{\Gal(K^{\mathrm{sep}}/K)}(V_\ell A)
\]
is an isomorphism.
\end{theorem}

\begin{proof}[Source proof]

Let $W\subset V_\ell A$ be Galois-stable and set
$\mathfrak a=\{u\in\End(A)\otimes\Q_\ell:uV_\ell A\subset W\}$.  This is a
right ideal, and the stable-subspace lemma gives
$\mathfrak aV_\ell A=W$.  Since $\End(A)\otimes\Q_\ell$ is semisimple,
$\mathfrak a=e(\End(A)\otimes\Q_\ell)$ for an idempotent $e$.  Therefore
$V_\ell A=eV_\ell A\oplus(1-e)V_\ell A$, and the second summand is stable
because Galois commutes with $\End(A)\otimes\Q_\ell$.

For (b), let $C$ be the centralizer of $\End(A)\otimes\Q_\ell$ in
$\End(V_\ell A)$ and $B$ its centralizer.  By the double-centralizer theorem,
$B=\End(A)\otimes\Q_\ell$.  If $\varphi\in\End(V_\ell A)$, its graph is a
Galois-stable subspace of $V_\ell(A\times A)$.  The stable-subspace lemma
produces $u\in\End(A\times A)\otimes\Q_\ell$ with image this graph.  For
$c\in C$, the diagonal endomorphism
$\left(\begin{smallmatrix}c&0\\0&c\end{smallmatrix}\right)$ commutes with $u$,
so it preserves the graph.  Hence $\varphi c=c\varphi$ for every $c\in C$,
and $\varphi\in B$ as required.

\end{proof}


\section{Finiteness I implies Finiteness II}

\begin{theorem}[Bounded extensions]
\label{thm:milne-bounded-number-fields}

\dcref{MI:ch4:3.1}
For any number field $K$, integer $N$, and finite set of primes $S$ of $K$,
only finitely many fields $L\supset K$, up to $K$-isomorphism, are unramified
outside $S$ and have degree $N$.
\end{theorem}

\begin{proof}[Source proof]

For each place $v$ and each degree dividing $N$, there are only finitely many
extensions of $K_v$ of that degree.  This follows from Krasner's lemma: such
extensions are represented by monic integral polynomials of bounded degree,
and the compactness of the coefficient space gives finitely many sufficiently
close equivalence classes.  If $L/K$ is unramified outside $S$, its relative
discriminant is the product of the local discriminants at places above $S$;
the local finiteness therefore bounds it.  Hermite's theorem gives only
finitely many global fields with bounded discriminant.

\end{proof}


\begin{theorem}[Neron criterion for good reduction]
\label{thm:milne-good-reduction-criterion}

\dcref{MI:ch4:3.5}
\uses{thm:milne-bounded-number-fields,thm:milne-neron-model}
For a finite place $v$ and a prime $\ell$ different from its residue
characteristic, an abelian variety has good reduction at $v$ exactly when the
representation on $V_\ell A$ is unramified at $v$.
\end{theorem}

\begin{proof}[Source proof]

Let $K_v^{\rm un}$ and $K_v^{\rm al}$ denote the maximal unramified and
algebraic extensions of $K_v$, and let $A^0$ be the special fibre of the
Neron model.  For $\ell\ne\operatorname{char}(k(v))$, reduction gives a
bijection

\[
A(K_v)[\ell^n]\;\xrightarrow{\sim}\;A^0(k(v))[\ell^n],
\]

and the same remains true after replacing $K_v$ by an unramified extension.
The algebraic group $A^0$ has a filtration whose successive quotients are a
finite algebraic group, an abelian variety $B$, a torus $T$, and a unipotent
group $U$, with
$g=\dim B+\dim T+\dim U$.  Over an algebraic closure their $\ell^n$-torsion
has respectively $\ell^{2n\dim B}$, $\ell^{n\dim T}$, and $1$ elements.
If $A$ has good reduction, $A^0$ is an abelian variety, so
\[
\#A(K_v^{\rm un})[\ell^n]=\ell^{2ng}
 =\#A(K_v^{\rm al})[\ell^n].
\]
Thus inertia acts trivially on every $\ell$-power torsion point, and the
representation on $V_\ell A$ is unramified.  Conversely, if $A$ does not
have good reduction, then for all sufficiently large $n$,
$\#A(K_v^{\rm un})[\ell^n]<\ell^{2ng}$.  But
$A(K_v^{\rm un})[\ell^n]$ is the inertia-invariant subgroup of
$A(K_v^{\rm al})[\ell^n]$, so inertia does not act trivially on
$V_\ell A$.

\end{proof}


\begin{theorem}[Finiteness I implies Finiteness II]
\label{thm:milne-finiteness-I-II}

\dcref{MI:ch4:3.10}
\uses{thm:milne-tate-semisimplicity,thm:milne-good-reduction-criterion,thm:milne-bounded-number-fields,thm:milne-finite-isogeny-possibilities}
Assume finiteness of $K$-isomorphism classes in each isogeny class.  Then, for
fixed dimension and finite $S$, only finitely many $K$-isomorphism classes of
abelian varieties have good reduction outside $S$.
\end{theorem}

\begin{proof}[Source proof]

Fix a prime $\ell$ away from the residue characteristics in $S$ and take the
finite set $T=T(S,\ell,g)$ from Theorem~\ref{thm:milne-finite-isogeny-possibilities}.
For an abelian variety $A/K$ of dimension $g$ with good reduction outside $S$,
each $P_v(A,t)$, $v\in T$, is a monic integral polynomial of degree $2g$ whose
roots have absolute value $N(v)^{1/2}$.  There are only finitely many such
polynomials for each $v$, hence only finitely many tuples
$(P_v(A,t))_{v\in T}$.  The finite-isogeny theorem shows that each tuple
determines at most one isogeny class, and Finiteness~I makes each isogeny class
contain only finitely many $K$-isomorphism classes.

\end{proof}


\begin{definition}[Frobenius polynomial at a good place]
\label{def:milne-frobenius-polynomial}

\dcref{MI:ch4:3.4}
Let $A/K$ have good reduction at a finite place $v$.  If $A_v$ is its special
fibre and $\pi_{A_v}$ is its Frobenius endomorphism, define
\[
P_v(A,t)=\det(t-\pi_{A_v}\mid V_\ell A_v).
\]
For $\ell$ different from the residue characteristic of $v$, this polynomial
is independent of $\ell$ and has integral coefficients.
\end{definition}

\section{Finiteness II and Shafarevich}

\begin{proposition}[Jacobians preserve good reduction]
\label{prop:milne-jacobian-good-reduction}

\dcref{MI:ch4:4.1}
\uses{thm:milne-jacobian-general,def:milne-polarization}
If a smooth projective curve has good reduction at a place, its Jacobian has
good reduction there and carries the canonical principal polarization.  The
construction is compatible with the integral model and with base change.
\end{proposition}

\begin{proof}[Source proof]

Let $R$ be the local ring of $K$ at the place $v$.  Good reduction of $C$
means that there is a proper smooth curve $\mathcal C/R$ with generic fibre
$C$ and special fibre $C_v$.  The relative Jacobian construction of
Chapter~III, applied over $R$, gives an abelian scheme $\mathcal J/R$.
Its generic fibre is $J(C)$ and its special fibre is $J(C_v)$, so $J(C)$ has
good reduction at $v$ and the reduction is the Jacobian of $C_v$.

\end{proof}


\begin{theorem}[Shafarevich finiteness for curves]
\label{thm:milne-shafarevich-curves}

\dcref{MI:ch4:4.3}
\uses{thm:milne-finiteness-I-II,prop:milne-jacobian-good-reduction,thm:milne-torelli-source}
Assume Finiteness II.  For fixed $K$, genus $g\ge2$, and finite $S$, only
finitely many $K$-isomorphism classes of complete nonsingular curves of genus
$g$ have good reduction outside $S$.
\end{theorem}

\begin{proof}[Source proof]

Finiteness~II may be stated for principally polarized abelian varieties, since
each underlying isomorphism class has only finitely many principal
polarizations.  By the preceding proposition, the canonical principally
polarized Jacobian of a curve with good reduction outside $S$ also has good
reduction outside $S$.  The rational Torelli theorem says that the map
$C\mapsto(J(C),\lambda(C))$ is injective on $K$-isomorphism classes.  Applying
Finiteness~II to these polarized Jacobians gives the required finiteness.

\end{proof}


\section{Shafarevich and Mordell}

\begin{theorem}[De Franchis finiteness]
\label{thm:milne-de-franchis}

\dcref{MI:ch4:5.2}
For curves $C'$ and $C$ over a field $k$, with $g(C)\ge2$, there are only
finitely many nonconstant maps $C'\to C$.
\end{theorem}

\begin{proof}[Source proof]

Let $g'=g(C')$ and let $d$ be the degree of a nonconstant map.  Hurwitz gives
$2g'-2=d(2g(C)-2)+\operatorname{ramification}$, so $d$ is bounded.  The graph
$\Gamma\subset X=C'\times C$ is isomorphic to $C'$, and
\[
\Gamma\cdot K_X=(2g(C)-2)d+(2g'-2).
\]
Adjunction gives $\Gamma\cdot K_X+\Gamma^2=2g'-2$, hence
$\Gamma^2=-(2g(C)-2)d$.  Thus the self-intersections of these graphs lie in a
finite set of negative integers.  For each Hilbert polynomial, the Hilbert
scheme of curves in $X$ is a finite union of varieties.  At a graph its
tangent space is $H^0(\Gamma,N_{\Gamma/X})$, which vanishes because
$\deg N_{\Gamma/X}=\Gamma^2<0$.  Each component containing such a graph is
therefore zero-dimensional, so there are only finitely many graphs and hence
only finitely many maps.

\end{proof}


\begin{theorem}[Parshin covering construction]
\label{thm:milne-parshin-cover}

\dcref{MI:ch4:5.1}
\uses{prop:milne-generalized-jacobian-coverings-source,thm:milne-de-franchis,thm:milne-bounded-number-fields,prop:milne-weak-mordell-weil-extension,lem:milne-parshin-cover-construction}
Let $K$ be a number field and let $S$ contain the dyadic primes.  For every
complete nonsingular curve $C/K$ of genus $g\ge1$ with good reduction outside
$S$, there is a finite extension $L/K$ such that for every $P\in C(K)$ there
are a curve $C_P/L$ and a finite map $\phi_P:C_P\to C_L$ satisfying:
\begin{enumerate}
\item $C_P$ has good reduction outside the places of $L$ above $S$;
\item the genus of $C_P$ is bounded independently of $P$;
\item $\phi_P$ is ramified exactly at $P$.
\end{enumerate}
\end{theorem}

\begin{proof}[Source proof]

If $C(K)$ is empty there is nothing to prove.  Otherwise embed $C$ in its
Jacobian $J$.  Multiplication by $2$ on $J$ is etale of degree $2^{2g}$; its
pullback gives an etale cover $C_1\to C$ of that degree.  Over the ring of
$S$-integers, the multiplication map is etale on the abelian scheme and the
embedding extends to the smooth proper models, so $C_1$ has good reduction
outside $S$.  The points above any $P\in C(K)$ lie in extensions unramified
outside $S$ and of degree at most $2^{2g}$.  The bounded-extension theorem
therefore gives one finite extension $L_1/K$, unramified outside $S$, over
which all these points are rational.

Choose two distinct points $P_1,P_2$ above $P$.  By the weak Mordell--Weil
proposition, after a further finite extension $L_2/L_1$ independent of $P$,
there is a divisor $D$ with $2D\sim P_1-P_2$.  Pass to a Hilbert class field
$L_3$ so that the relevant ring of integers is principal, and choose
$f\in L_3(C)^\times$ with
$(f)=P_1-P_2+2D$.  The square-root cover from
Lemma~\ref{lem:milne-parshin-cover-construction} is ramified exactly over
$P_1$ and $P_2$; composed with $C_1\to C$, it gives a map
$C_P\to C_{L_3}$ ramified exactly over $P$.  The Hurwitz formula bounds
$g(C_P)$ in terms of $g$ and the fixed degree $2^{2g+1}$, independently of
$P$.  The same lemma gives good reduction outside the places above $S$.

\end{proof}


\section{Numbered finiteness ingredients}

\begin{lemma}[Finite-level Tate sequences]
\label{lem:milne-tate-finite-level-sequences}

\dcref{MI:ch4:2.2}
\uses{lem:milne-finite-level-tate}
For $\ell\ne\operatorname{char}(K)$ and $n\ge1$,
\[
0\longrightarrow T_\ell A\xrightarrow{\ell^n}T_\ell A
\longrightarrow A[\ell^n](K^{\mathrm{sep}})\longrightarrow0.
\]
If $\alpha:A\to B$ is an isogeny of degree prime to
$\operatorname{char}(K)$, then
\[
0\longrightarrow T_\ell A\xrightarrow{T_\ell\alpha}T_\ell B
\longrightarrow C\longrightarrow0,
\]
where $C$ is finite of order equal to the $\ell$-primary part of
$\deg(\alpha)$.
\end{lemma}

\begin{proof}

The first sequence follows directly from
$T_\ell A=\{(a_n)_{n\ge1}:a_n\in A[\ell^n](K^{\mathrm{sep}}),\ \ell a_n=a_{n-1}\}$.
For the second, let $K_n$ and $C_n$ be the kernel and cokernel in the diagram
of finite torsion groups induced by $\alpha$.  For sufficiently large $n$,
$K_n=K_{n+1}=\cdots=K$.  Since $K$ is finite, it has no nonzero element
divisible by every power of $\ell$, so $\varprojlim K_n=0$.  The two torsion
groups have the same order; hence $\#K_n=\#C_n$, and $\#C_n$ is constant for
large $n$.  The transition maps $C_{n+1}\to C_n$ are then surjective and
therefore bijective.  Passing to inverse limits gives the displayed exact
sequence, and its finite cokernel has order the $\ell$-primary part of
$\deg(\alpha)$.

\end{proof}


\begin{lemma}[Stable lattices and isogenies]
\label{lem:milne-stable-tate-lattices}

\dcref{MI:ch4:2.3}
\uses{lem:milne-tate-finite-level-sequences}
Every finite-index Galois-stable sublattice of $T_\ell A$ is the image of an
$\ell$-power isogeny $A'\to A$ defined over $K$.  Conversely, the Tate module
of an $\ell$-power isogeny is such a stable sublattice.
\end{lemma}

\begin{proof}

Choose $n$ with $\ell^nT_\ell A\subset W$, and let $N$ be the image of $W$ in
$T_\ell A/\ell^nT_\ell A=A[\ell^n](K^{\mathrm{sep}})$.  By Galois stability,
$N$ is stable, so $B=A/N$ and the quotient map $\pi:A\to B$ are defined over
$K$.  The map $\ell^n:A\to A$ factors through $\pi$, say
$\ell^n=\beta\circ\pi$ with $\beta:B\to A$.  Consequently the image of
$T_\ell\beta$ contains $\ell^nT_\ell A$.  Its image modulo
$\ell^nT_\ell A$ is $N$: indeed
$B[\ell^n]=\{a\in A(K^{\mathrm{sep}}):\ell^na\in N\}/N$, and $\beta$ maps
this group onto $N$.  Thus $\operatorname{im}(T_\ell\beta)=W$, as required.

\end{proof}


\begin{corollary}[Tate double-centralizer]
\label{cor:milne-tate-double-centralizer}

\dcref{MI:ch4:2.7}
\uses{cor:milne-tate-hom-finiteness,thm:milne-tate-semisimplicity}
Under the Finiteness I hypotheses for the varieties occurring in the product,
the image $R$ of $\Q_\ell[\Gal(K^{\mathrm{sep}}/K)]$ in
$\End_{\Q_\ell}(V_\ell A)$ is the centralizer of
$\End^0(A)\otimes\Q_\ell$.
\end{corollary}

\begin{proof}[Source proof]

Theorem~2.5 makes $V_\ell A$ a faithful semisimple $R$-module.  The double
centralizer theorem, together with the endomorphism equality in Theorem~2.5,
identifies the centralizer of $\End^0(A)\otimes\Q_\ell$ with $R$.

\end{proof}


\begin{theorem}[Hermite finiteness]
\label{thm:milne-hermite-finiteness}

\dcref{MI:ch4:3.2}
\uses{thm:milne-bounded-number-fields}
For any number field $K$, integer $N$, and finite set of primes $S$ of $K$,
there are only finitely many fields $L\supset K$, up to $K$-isomorphism, that
are unramified outside $S$ and have degree $N$.  Equivalently, there are only
finitely many number fields of a fixed discriminant.
\end{theorem}

\begin{proof}[Source proof]

Let $D=|\operatorname{Disc}(L/\mathbb Q)|$.  Minkowski's ideal-class bound
gives an integral ideal in every ideal class whose norm is bounded in terms of
$[L:\mathbb Q]$ and $D$.  Since such a norm is at least one, a bound on $D$
also bounds $[L:\mathbb Q]$ (the factorial term in Minkowski's bound dominates
as the degree grows).  For a fixed degree, embed $L$ in its real and complex
Minkowski space.  Minkowski's theorem supplies a nonzero algebraic integer
whose conjugates are bounded in terms of $D$.  The coefficients of its minimal
polynomial are consequently bounded, so only finitely many such polynomials
can occur; the usual choice in the Minkowski argument generates $L$ over
$\mathbb Q$.  Thus only finitely many fields have the given discriminant, and
the bounded-extension assertion follows.

\end{proof}


\begin{theorem}[Chebotarev density]
\label{thm:milne-chebotarev-detection}

\dcref{MI:ch4:3.3}
Let $L/K$ be a finite Galois extension of number fields with group $G$, and let
$C\subset G$ be stable under conjugation.  The primes $v$ of $K$ with
Frobenius conjugacy class $(v,L/K)=C$ have density $|C|/|G|$.
\end{theorem}

\begin{proof}[Source-omitted proof]

This is the Chebotarev density theorem; Milne cites ANT, 8.31, and CFT,
VIII~\S7 for the proof.

\end{proof}


\begin{corollary}[Good reduction is preserved by isogeny]
\label{cor:milne-isogeny-good-reduction}

\dcref{MI:ch4:3.6}
\uses{thm:milne-good-reduction-criterion}
If $A$ and $B$ are isogenous over $K$, and one has good reduction at $v$, then
the other has good reduction at $v$.
\end{corollary}

\begin{proof}[Source proof]

The isogeny induces an isomorphism $V_\ell A\simeq V_\ell B$ compatible with
the Galois action.  The Neron criterion (Theorem~3.5) gives the result.

\end{proof}


\begin{theorem}[Finite isogeny possibilities]
\label{thm:milne-finite-isogeny-possibilities}

\dcref{MI:ch4:3.7}
\uses{thm:milne-good-reduction-criterion,thm:milne-tate-semisimplicity,cor:milne-tate-hom-finiteness,lem:milne-semisimple-module-finiteness,thm:milne-chebotarev-detection,def:milne-frobenius-polynomial}
Let $A,B$ be abelian varieties of dimension $g$ over $K$.  Let $S$ contain all
places where either has bad reduction, and choose $\ell$ away from the residue
characteristics in $S$.  There is a finite set
$T=T(S,\ell,g)$, disjoint from $S$ and the places above $\ell$, such that
\[
 P_v(A,t)=P_v(B,t)\quad(v\in T)\quad\Longrightarrow\quad A\text{ and }B
\text{ are isogenous}.
\]
\end{theorem}

\begin{proof}[Source proof]

By the Neron criterion and Tate semisimplicity, $V_\ell A$ and $V_\ell B$ are
semisimple representations unramified outside $S\cup\{v:v\mid\ell\}$.  Apply
the semisimple-module finiteness lemma to obtain a finite set $T$ of places.
At a place outside $S$ and above $\ell$, the Frobenius characteristic
polynomial on $V_\ell A$ is $P_v(A,t)$, and similarly for $B$.  Equality at
all $v\in T$ therefore gives an isomorphism
$V_\ell A\simeq V_\ell B$ of Galois representations.  Tate's homomorphism
theorem then gives an isogeny $A\to B$.

\end{proof}


\begin{lemma}[Semisimple module finiteness]
\label{lem:milne-semisimple-module-finiteness}

\dcref{MI:ch4:3.8}
\uses{thm:milne-tate-semisimplicity}
Let $(V,\rho)$ and $(W,\sigma)$ be semisimple representations of
$\Gal(K^{\mathrm{sep}}/K)$ on $\Qell$-vector spaces of dimension $d$, both
unramified outside $S\cup\{v:v\mid\ell\}$.  There is a finite set
$T=T(S,\ell,d)$, disjoint from that bad set, such that equality of the
characteristic polynomials of Frobenius on $V$ and $W$ at every $v\in T$ implies
$(V,\rho)\simeq(W,\sigma)$.
\end{lemma}

\begin{proof}[Source proof]

By Hermite finiteness, the composite $L$ of all subextensions of
$K^{\mathrm{sep}}/K$ of degree at most $\ell^{2d}$ unramified outside
$S\cup\{v:v\mid\ell\}$ is finite, Galois, and unramified outside that set.
Chebotarev supplies a finite set $T$ of places, disjoint from the bad set,
whose Frobenius conjugacy classes cover $\operatorname{Gal}(L/K)$.

Choose full stable $\mathbb Z_\ell$-lattices $M\subset V$ and $N\subset W$.
The $\mathbb Z_\ell$-algebra $R$ generated by the two actions of
$\operatorname{Gal}(L/K)$ acts on both lattices.  Equality of the Frobenius
characteristic polynomials gives equality of traces on the generators of $R$,
and hence on all of $R$.  The trace criterion for semisimple $R\otimes\Q_\ell$-
modules now gives $(V,\rho)\simeq(W,\sigma)$.

\end{proof}


\begin{lemma}[Trace criterion for semisimple modules]
\label{lem:milne-endomorphism-algebra-bound}

\dcref{MI:ch4:3.9}
\uses{thm:milne-tate-semisimplicity}
Let $k$ have characteristic zero and let $R$ be a finite-dimensional
$k$-algebra.  Two finite-dimensional semisimple $R$-modules are isomorphic if
and only if they have the same trace on every element of $R$.
\end{lemma}

\begin{proof}[Source-omitted proof]

This is the standard character criterion for semisimple modules (Milne cites
Bourbaki, \emph{Algebre}, VIII, \S12, no.~1, Prop.~3).

\end{proof}


\begin{theorem}[Rational Torelli for the Shafarevich step]
\label{thm:milne-rational-torelli}

\dcref{MI:ch4:4.2}
\uses{thm:milne-torelli-source,prop:milne-jacobian-good-reduction}
Let $C$ be a complete nonsingular curve of genus at least $2$ over a perfect
field $k$.  The isomorphism class of $C$ is uniquely determined by the
principally polarized abelian variety $(J(C),\lambda(C))$.
\end{theorem}

\begin{proof}[Source proof]

After extending to an algebraic closure, the Torelli theorem of
Chapter~III gives an isomorphism $\alpha:C_{k^{\rm al}}\to C'_{k^{\rm al}}$
from an isomorphism of the polarized Jacobians.  The theorem also specifies
$\alpha$ uniquely once the relevant points and sign are fixed (and gives the
same curve isomorphism independently of the auxiliary points in the rational
statement).  For every $\sigma\in\operatorname{Gal}(k^{\rm al}/k)$, the
isomorphism associated with $\sigma\beta$ is $\sigma\alpha$.  Since $\beta$ is
defined over $k$, $\sigma\beta=\beta$, and uniqueness gives
$\sigma\alpha=\alpha$.  Thus $\alpha$ descends to $k$; perfectness makes this
the required effective descent of curves.

\end{proof}


\begin{lemma}[Square-root cover construction]
\label{lem:milne-parshin-cover-construction}

\dcref{MI:ch4:5.3}
\uses{thm:milne-parshin-cover}
Let $C/K$ be complete nonsingular and let
$(f)=P_1-P_2+2D$ with $P_1,P_2\in C(K)$.  The double cover obtained from
$K(C)\subset K(C)(\sqrt f)$ is ramified exactly at $P_1$ and $P_2$.  If $S$
contains the bad-reduction primes of $C$, the primes where $P_1$ and $P_2$
coalesce, and the primes above $2$, and the ring of $S$-integers is principal,
then the cover has good reduction outside $S$.
\end{lemma}

\begin{proof}[Source proof]

The odd valuations of $f$ are exactly those at $P_1$ and $P_2$.  Extend $C$
to a smooth model over the $S$-integers; principality lets one multiply $f$ by
a constant so that its divisor has no vertical components.  The resulting
double cover is smooth away from the dyadic and collision primes.

\end{proof}


\begin{proposition}[Weak Mordell--Weil field extension]
\label{prop:milne-weak-mordell-weil-extension}

\dcref{MI:ch4:5.5}
\uses{thm:milne-parshin-cover}
For an abelian variety $A/K$ with good reduction outside a finite set, there
is a finite extension $L/K$ such that every point of $A(K)$ is twice a point
of $A(L)$.
\end{proposition}

\begin{proof}[Source proof]

The quotient $A(K)/2A(K)$ is finite by Mordell--Weil.  Adjoin to $K$ the
coordinates of representatives of this finite quotient.

\end{proof}


\begin{proposition}[Elliptic height normalization choices]
\label{prop:milne-faltings-height-normalization}

\dcref{MI:ch4:6.1}
\uses{def:milne-faltings-height}
For a complex elliptic curve $E$, each of the following choices determines the
remainder:
\begin{enumerate}
\item an isomorphism $\C/\Lambda\to E(\C)$;
\item a basis of $(E,\Omega^1)$, equivalently a nonzero holomorphic
differential;
\item an equation $Y^2=4X^3-g_2X-g_3$ for $E$.
\end{enumerate}
\end{proposition}

\begin{proof}[Source proof]

The Weierstrass function attached to $\Lambda$ gives the equation and the
uniformization.  The equation gives $dX/Y$, while integrating a nonzero
differential over a homology basis recovers the lattice.

\end{proof}


\begin{lemma}[Height of a normed rank-one module]
\label{lem:milne-faltings-height-formula}

\dcref{MI:ch4:6.2}
\uses{prop:milne-faltings-height-normalization}
The height defined from a normed rank-one module is independent of the chosen
nonzero module element.
\end{lemma}

\begin{proof}[Source proof]

Write the finite-index factors using local generators.  Their product is the
product of normalized absolute values of the scalar relating two generators,
which is $1$ by the product formula.

\end{proof}


\begin{proposition}[Invariant differentials]
\label{prop:milne-arakelov-determinant}

\dcref{MI:ch4:6.4}
\uses{def:milne-faltings-height}
If $V$ is a smooth variety of dimension $g$, then $\Omega^1_{V/k}$ is locally
free of rank $g$; when $V$ is a group variety it is free.
\end{proposition}

\begin{proof}[Source proof]

This is the standard smoothness criterion for differentials; for a group
variety, translation identifies the sheaf with its fibre at the identity.

\end{proof}


\begin{corollary}[Top differential sheaf]
\label{cor:milne-height-finite-maps}

\dcref{MI:ch4:6.5}
\uses{prop:milne-arakelov-determinant}
The sheaf $\Omega^g_{V/k}=\bigwedge^g\Omega^1_{V/k}$ is locally free of rank
one, and it is free when $V$ is a group variety.
\end{corollary}

\begin{proof}[Source proof]

Take the top exterior power of the locally free sheaf in Proposition~6.4.

\end{proof}


\begin{proposition}[Global sections on a complete variety]
\label{prop:milne-bounded-complete-height}

\dcref{MI:ch4:6.6}
\uses{cor:milne-height-finite-maps}
If $V$ is complete and geometrically connected and $M$ is a free coherent
sheaf, evaluation at any $k$-rational point identifies $H^0(V,M)$ with the
fibre $M(v)$.
\end{proposition}

\begin{proof}[Source proof]

Complete connected varieties have only constant global functions.  For a free
sheaf $M\simeq\mathcal O_V^{\oplus n}$, evaluation is therefore $k^n\to k^n$.

\end{proof}


\begin{proposition}[Invariant differential evaluation]
\label{prop:milne-isogeny-height-estimate}

\dcref{MI:ch4:6.7}
\uses{cor:milne-height-finite-maps}
For an abelian variety $A$ of dimension $g$, evaluation at the identity gives
canonical isomorphisms $H^0(A,\Omega^1_{A/k})\simeq\Omega^1_{A/k}(0)$ and
$H^0(A,\Omega^g_{A/k})\simeq\Omega^g_{A/k}(0)$.
\end{proposition}

\begin{proof}[Source proof]

Combine the preceding two propositions and evaluate at the identity element.

\end{proof}


\begin{theorem}[Siegel moduli over the complex numbers]
\label{thm:milne-siegel-moduli-complex}

\dcref{MI:ch4:7.3}
\uses{def:milne-polarization}
For fixed dimension $g$ and polarization degree $d$, there is a unique
algebraic variety $M_{g,d}/\C$ and a bijection between $M_{g,d}(\C)$ and
isomorphism classes of pairs $(A,\lambda)$ of dimension $g$ and degree $d$,
such that locally every point is represented by a family and every family over
a variety induces a regular classifying map to $M_{g,d}$.
\end{theorem}

\begin{proof}[Source proof]

Uniqueness follows by applying the family property in neighborhoods and gluing
the resulting inverse regular maps.  Existence is the algebraicity theorem for
the Siegel modular variety.

\end{proof}


\begin{proposition}[Northcott finiteness in projective space]
\label{prop:milne-projective-height}

\dcref{MI:ch4:7.1}
For a number field $K$ and a constant $C$, only finitely many points of
projective space over $K$ have absolute Weil height at most $C$.
\end{proposition}

\begin{proof}[Source proof]

Clear denominators and use the product formula to choose primitive integral
coordinates.  A height bound bounds every coordinate, leaving finitely many
tuples.

\end{proof}


\begin{proposition}[Height comparison for projective embeddings]
\label{prop:milne-projective-height-comparison}

\dcref{MI:ch4:7.2}
\uses{prop:milne-projective-height}
If two projective embeddings of a variety pull back linearly equivalent
hyperplanes, their associated Weil height functions differ by a bounded amount.
\end{proposition}

\begin{proof}[Source proof]

The ratio of the two height functions is represented by a rational function
with divisor zero, hence its logarithmic absolute values are bounded on the
projective variety.

\end{proof}


\begin{theorem}[Canonical modular model]
\label{thm:milne-siegel-moduli-rational-model}

\dcref{MI:ch4:7.5}
\uses{thm:milne-siegel-moduli-complex}
There is a unique model of $M_{g,d}$ over $\Q$ for which the complex
classification bijection commutes with the two actions of $\operatorname{Aut}(\C)$.
For every field $L\supset\Q$ this gives a functorial map
$M_{g,d}(L)\to\{(A,\lambda)\text{ over }L\}/\cong$, which is an isomorphism
when $L$ is algebraically closed.
\end{theorem}

\begin{proof}[Source proof]

The complex classification is equivariant for $\operatorname{Aut}(\C)$, so it
descends uniquely to $\Q$.  Base change gives the functorial map, and over an
algebraically closed field every polarized pair has a complex classification.

\end{proof}


\begin{theorem}[Faltings and modular height comparison]
\label{thm:milne-bounded-modular-height}

\dcref{MI:ch4:7.8}
\uses{thm:milne-siegel-moduli-rational-model,prop:milne-projective-height}
For every polarized abelian variety over a number field,
$h_F(A)=h_M(A,\lambda)+O(\log h_M(A,\lambda))$.
\end{theorem}

\begin{proof}[Source-omitted proof]

Milne explains that this comparison is obtained by extending the two height
functions to a compactification of the Siegel modular scheme over $\mathbb Z$;
the analysis of their boundary singularities is not proved in the source.

\end{proof}


\begin{theorem}[Modular finiteness with fixed height]
\label{thm:milne-modular-finiteness-fixed-level}

\dcref{MI:ch4:7.10}
\uses{thm:milne-bounded-modular-height,prop:milne-projective-height}
Let $K$ be a number field and $g,d,C$ integers.  Up to isomorphism there are
only finitely many polarized abelian varieties $(A,\lambda)$ over $K$ of
dimension $g$ and degree $d$, semistable everywhere, with $h_M(A,\lambda)\le C$.
\end{theorem}

\begin{proof}[Source proof]

The finite set of points of $M_{g,d}(K)$ of bounded height is supplied by
Northcott.  Fix one geometric pair $(A_0,\lambda_0)$ among them and choose
$\ell\ge3$.  The field generated by the points of order $\ell$ on a pair in
this geometric class has degree at most
$\#\operatorname{GL}_{2g}(\mathbb F_\ell)$ and is unramified outside the
bad places and the places above $\ell$.  Hermite finiteness therefore gives a
single finite Galois extension $L/K$ over which the $\ell$-torsion of every
such pair is rational.

If $\alpha:(A,\lambda)\to(A_0,\lambda_0)$ is an isomorphism over
$K^{\mathrm{sep}}$ and $\sigma\in\operatorname{Gal}(K^{\mathrm{sep}}/L)$,
then $\sigma\alpha$ and $\alpha$ agree on $A[\ell]$.  Their quotient is an
automorphism of the polarized pair $(A_0,\lambda_0)$ acting trivially on
$A_0[\ell]$, hence is the identity.  Thus $\alpha$ is defined over $L$.
After replacing $L$ by a Galois extension, the $K$-forms are classified by
the finite set
$H^1(\operatorname{Gal}(L/K),\operatorname{Aut}(A_{0,L},\lambda_{0,L}))$.
This proves finiteness.

\end{proof}


\begin{corollary}[Faltings-height polarized finiteness]
\label{cor:milne-unpolarized-height-consequence}

\dcref{MI:ch4:7.12}
\uses{thm:milne-modular-finiteness-fixed-level,thm:milne-bounded-modular-height}
For fixed dimension, polarization degree, number field, and semistability,
there are only finitely many polarized abelian varieties of bounded Faltings
height.
\end{corollary}

\begin{proof}[Source proof]

The height comparison theorem turns a Faltings-height bound into a modular-height
bound, so apply Theorem~7.10.

\end{proof}


\begin{corollary}[Semistable unpolarized finiteness]
\label{cor:milne-semistable-unpolarized-finiteness}

\dcref{MI:ch4:7.13}
\uses{cor:milne-unpolarized-height-consequence}
For fixed $K$, dimension $g$, and real $C$, there are only finitely many
abelian varieties over $K$ of dimension $g$ with semistable reduction everywhere
and $h_F(A)\le C$, up to isomorphism.
\end{corollary}

\begin{proof}[Source proof]

Use the principally polarized abelian variety $(A\times A^\vee)^4$ and the
identity $h_F(A^\vee)=h_F(A)$.  Its Faltings height is $8h_F(A)$, so the
polarized finiteness result applies.

\end{proof}


\section{Faltings height}

\begin{definition}[Faltings height]
\label{def:milne-faltings-height}

\dcref{MI:ch4:6}
For a semistable abelian variety over a number field, the stable Faltings height
is defined from the Arakelov degree of the determinant of invariant
differentials, with the archimedean metric induced by a period integral.  The
definition is normalized by an auxiliary multiplication integer and then shown
to be independent of that choice.
\end{definition}

\begin{lemma}[Height of a normed module under field extension]
\label{lem:milne-faltings-height-isogeny}

\dcref{MI:ch4:6.3}
\uses{def:milne-faltings-height,def:milne-polarization}
Let $M$ be a normed projective rank-one module over the integers of $K$.  In
the expression for $H(M)$ a nonzero element of $M\otimes K$ may be used, and
for $h(M)=[K:\Q]^{-1}\log H(M)$ one has
\[
h(\mathcal O_L\otimes_{\mathcal O_K}M)=h(M)
\]
for every finite extension $L/K$.
\end{lemma}
\begin{proof}
Choose a nonzero $m\in M\otimes K$.  The determinant norm in the extension
is the field norm of the original determinant, so its logarithm is multiplied
by $[L:K]$.  Dividing by $[L:\Q]=[L:K][K:\Q]$ gives the same normalized
height.
\end{proof}

\section{Completion of Finiteness I}

\begin{theorem}[Faltings-height finiteness in an isogeny class]
\label{thm:milne-complete-finiteness-I}

\dcref{MI:ch4:8.1}
\uses{lem:milne-faltings-height-isogeny,thm:milne-bounded-modular-height,thm:milne-modular-finiteness-fixed-level}
Let $A/K$ have semistable reduction everywhere.  The set of stable Faltings
heights of abelian varieties $B/K$ isogenous to $A$ is finite.
\end{theorem}

\begin{proof}[Source-omitted proof]

Milne states that Faltings proves the height assertion using algebraic
geometry, including Raynaud's finiteness theorem for finite group schemes, but
does not give the argument.  Together with the polarized bounded-height
finiteness theorem, it yields Finiteness~I.

\end{proof}
